\documentclass[10pt,a4paper]{article}
\usepackage{fontspec,amsmath,amssymb,graphicx}
\usepackage[margin=18mm]{geometry}
\setmainfont{Times New Roman}
\setlength{\parindent}{0pt}
\setlength{\parskip}{5pt}
\allowdisplaybreaks
\sloppy
\emergencystretch=2em
\begin{document}
\section*{BT\_5\_5\_3\_08 — Lời giải bổ sung}
Nhân Poisson dương nên $h\ge0$ suy $u\ge0$. Nếu $h$ không đồng nhất 0 và liên tục, tích phân nhân dương cho $u>0$. Nếu $h=0$, duy nhất Dirichlet cho $u=0$.
Với $r=1/2$, $P_r(\theta)=(1-r^2)/(1-2r\cos\theta+r^2)$ nằm giữa $(1-r)/(1+r)=1/3$ và $(1+r)/(1-r)=3$. Vì trung bình $h$ bằng $u(0)=1$, tích phân cho $1/3\le u\le3$.\end{document}